% NMLT Paper 1 -- post-pivot draft
% Normative Lean:
%   mechanization/lean/NMLT/Behavior/ResourceBehavior.lean
%   mechanization/lean/NMLT/Artifact/BehaviorCore.lean
%   mechanization/lean/NMLT/Artifact/SemanticClosure.lean
% Normative source:
%   examples/pivot/visible_resource_sync.nmlt
\documentclass[11pt]{article}

\usepackage[margin=1in]{geometry}
\usepackage{amsmath,amssymb,amsthm}
\usepackage{booktabs}
\usepackage{microtype}
\usepackage[hyphens]{url}
\usepackage[hidelinks]{hyperref}

\newtheorem{theorem}{Theorem}
\newtheorem{definition}[theorem]{Definition}
\newtheorem{remark}[theorem]{Remark}

\newcommand{\nmlt}{\textsc{Nmlt}}
\newcommand{\core}{\texttt{behavior-core-v1}}

\title{Resource-Aware Weak Refinement for an Open Programming Language\\
\large A Lean-defined behavioral core and a checked composition theorem}
\author{Carlosian\\\texttt{carlosian@agentmail.to}}
\date{30 August 2026}

\begin{document}
\maketitle

\begin{abstract}
\nmlt{} is a pre-alpha programming language for expressing finite open
behaviors with typed ports, affine capabilities, additive grades, and
rely/guarantee contracts.  Its normative behavioral semantics is defined in
Lean~4; Rust supplies a lossless frontend, typed elaboration, a deterministic
semantic artifact, and a non-authoritative reference explorer.

We define resource-aware weak refinement directly over the step relation of a
single behavior object and prove that it lifts through binary composition when
the wiring is preserved and both products are well formed.  Product formation
requires complementary port directions and equal payloads, excludes connected
hidden actions, partitions affine authority, matches transfer and receive
profiles exactly, and discharges synchronized relies with peer guarantees.
Refinement preserves capability effects, is pointwise monotone in additive
grades, and requires every hidden step to refine the resource-free stutter
profile.  The theorem is instantiated by one source program that transfers a
nominal capability exactly once while accumulating a nonzero grade.

This result is safety-oriented.  It does not state or transport fairness,
infinite-trace, divergence, or liveness properties, and it does not constitute
a verified Rust-to-Lean compiler.
\end{abstract}

\section{Language-to-mathematics boundary}

The first vertical slice has the following auditable path:
\[
 \texttt{.nmlt}
 \longrightarrow
 \text{resolved typed core}
 \longrightarrow
 \core
 \longrightarrow
 \text{Lean validation and theorem instantiation}.
\]
The checked source fixture is
\url{../../examples/pivot/visible_resource_sync.nmlt}; its byte-canonical
artifact is
\url{../../examples/pivot/visible_resource_sync.behavior-core-v1.json}.
The artifact records the source SHA-256 digest, typed term trees, ports,
transitions, complete resource profiles, compositions, connections,
refinement maps, and hidden-action maps.

Rust may reject a malformed program early, but Rust acceptance is not the
semantic authority.  The Lean decoder rejects stale digests, noncanonical
artifacts, duplicate identities and resource entries, malformed finite types,
invalid connections, unmatched authority transfer, undischarged relies, and
invalid refinement witnesses.  Exact snapshots bind the compiler output to a
source revision; they are an audit mechanism, not a compiler-correctness
theorem.

\section{Normative behavior}

\begin{definition}[Open resource behavior]
A behavior contains a state type, initial-state predicate, transition
relation, observation function, action visibility, action direction and
payload, nominal capability ownership, and a resource profile for every
action.  A resource profile contains required, consumed, transferred, and
received capabilities; a finite-map grade; relied-upon facts; and guaranteed
facts.
\end{definition}

Capabilities are affine nominal identities.  Grades map named atoms to natural
numbers.  Sequential and parallel combination use pointwise addition, zero is
the stutter grade, and a concrete grade refines an abstract grade pointwise.
Contract facts are constructors of a finite enum.  Parallel resource profiles
retain an undisclosed reliance precisely when the peer does not guarantee it.
Synchronization therefore requires mutual rely discharge.

Hiding changes observability only.  It does not erase capability effects,
grade, or reliance.  Consequently a hidden concrete transition may refine
abstract stutter only when the mapped state is unchanged and its complete
resource profile refines the empty profile.

\section{Binary composition}

The product step relation has three constructors: an isolated left step, an
isolated right step, and a synchronized step.  A connected action cannot step
alone.  A synchronized step exists only when both component transitions
exist.  Its resource profile is the parallel combination of the two actual
action profiles, so the theorem cannot substitute a detached
label-to-resource interpretation.

A formed product establishes:
\begin{enumerate}
  \item complementary input/output directions and equal typed payloads;
  \item no connected hidden action;
  \item disjoint component ownership;
  \item exact transfer/receive agreement in both directions; and
  \item discharge of every synchronized reliance by the peer guarantee.
\end{enumerate}
Wiring equivalence preserves and reflects the complete binary connection
relation.  It therefore preserves isolation as well as synchronization.

\section{Resource-aware weak refinement}

A witness maps concrete states to abstract states and observations to
observations.  It maps initial states to initial states, preserves the hidden
classification, maps hidden steps to state equality, and maps visible steps to
actual abstract steps.  For every action it also supplies resource refinement:
requirements do not grow, authority effects agree exactly, concrete grades do
not exceed abstract grades, concrete relies do not grow, and concrete
guarantees cover abstract guarantees.

\begin{theorem}[Conditional binary-composition lift]
Let a concrete behavior resource-weakly refine an abstract behavior.  Let a
peer be composed with each behavior using whole-wiring-equivalent
connections.  If both products are formed under the port, visibility,
ownership, exact-transfer, grade, and rely/guarantee rules above, then the
component state map paired with the identity on the peer is a resource-aware
weak refinement of the concrete product by the abstract product.
\end{theorem}

The checked declaration is
\texttt{NMLT.Behavior.ResourceBehavior.liftParallel}.  Its proof performs case
analysis on the actual product transition relation.  Isolated component steps
use the local refinement and reflected isolation; peer steps are preserved;
synchronized steps use wiring preservation and the visible-step clause.
Resource refinement of synchronized actions follows from pointwise grade
addition and parallel resource refinement.

\section{Source-level instantiation}

The primary fixture declares an abstract sender, a concrete sender, and a
receiver.  Both senders own \texttt{permit} and emit it over an output port
with payload \texttt{Once<Unit>}; the receiver binds that capability on the
matching input.  Sender work grades are two and one, respectively, while
receiver work is two.  The sender relies on \texttt{Ready} and guarantees
\texttt{Authorized}; the receiver supplies the dual contract.

The explicit state refinement is identity.  The concrete visible
synchronization moves authority from the sender to the receiver, changes the
  receiver's observed bit, and has composed work grade three.  Lean decodes the
  canonical artifact into the actual finite behaviors, checks the complete
  premise bundle, and applies
  \texttt{NMLT.Artifact.SemanticClosure.Certificate.lifted}; there is no
  hand-mirrored source example in the proof tree.

\section{Negative controls}

Permanent controls show that the theorem cannot simply drop each of the
following obligations:
\begin{center}
\begin{tabular}{ll}
\toprule
Omitted obligation & Witnessed failure\\
\midrule
hidden-boundary isolation & a hidden action remains connected\\
whole-wiring preservation & an isolated step becomes synchronized\\
capability partition & both components own the same authority\\
transfer matching & transfer lacks a matching receive\\
grade preservation & concrete work exceeds abstract work\\
rely discharge & the peer omits a required guarantee\\
hidden-step resources & hidden consumption is treated as stutter\\
\bottomrule
\end{tabular}
\end{center}
These declarations live in
\texttt{NMLT.Counterexamples.ResourceBehaviorControls}.  Source fixtures add
separate diagnostics for hidden consumption, hidden nonzero grade, hidden
extra reliance, incompatible ports, undischarged reliance, shared ownership,
unmatched transfer, and incomplete state maps.  None is accepted merely by a
generic hidden-action rejection.

\section{Trusted base and claim ceiling}

Lean~4.30 is the normative proof environment.  The active audit rejects
\texttt{sorry}, \texttt{admit}, project axioms, and unchecked native
decisions.  The two composition declarations above depend only on Lean's
standard propositional extensionality axiom, \texttt{propext}.  The retained
typed-elaboration metatheory additionally uses the documented
\texttt{Quot.sound}.

The Rust explorer is useful for bounded operational inspection, but it never
returns \emph{proved}, \emph{model checked}, or evidence claims.  No claim is
made here about liveness, fairness, infinite traces, probabilistic or hybrid
behavior, general composition, arbitrary grade algebras, code generation,
runtime attestation, or verified compilation.

\section{Next mathematics}

The next milestone may revisit hidden divergence only after defining
behavior-indexed fairness over the unified transition semantics.  Such a
development must distinguish infinite internal activity from unfair
scheduling and must state which fairness assumptions survive refinement and
composition.  Nothing in the present theorem transports those properties.

\end{document}
